\begin{pastexam}[label={px:nchu-114-dsa-20}]{中興資工所甲組 114 學年度 單選題第 20 題}
（單選；每題 3 分，未答題得 0 分，答錯不倒扣）

Which data structure is commonly used to implement breadth-first search (BFS) on graphs?

(A)~Stack \qquad (B)~Queue \qquad (C)~Heap \qquad (D)~Hash Table
\end{pastexam}

\begin{pastsoln}
答案是 (B)。

BFS 由起點一層一層往外擴：先處理距離 $0$ 的、再處理距離 $1$ 的、……。 用佇列，先發現的頂點先被拿出來
處理，佇列裡的頂點距離永遠是 $d$ 或 $d + 1$，所以依距離由近到遠處理（資結書第~21.3~節、虛擬碼~21.1）。
\begin{itemize}
\item (A) Stack：後進先出，會沿著最新發現的頂點一路往深處走，那是 DFS 的走法。 錯。
\item (C) Heap：依優先權取出，用在 Dijkstra、Prim 這類每次取最小值的演算法。 錯。
\item (D) Hash Table：可以用來記錄哪些頂點拜訪過（頂點不是 $1..n$ 的整數時），但決定拜訪順序的是佇列。 錯。
\end{itemize}
\end{pastsoln}
